%!TEX root = ../thesis.tex
\begin{cabstract}
随着互联网与移动终端的发展，图像、文本等多媒体数据持续增长，基于哈希的近似最近邻检索因存储开销低、检索速度快而被广泛应用于大规模检索任务。然而，多数现有方法建立在静态数据假设之上，难以适应真实场景中持续到达的数据流及其伴随的新类别出现、既有类别分布变化等概念漂移现象。概念漂移会破坏哈希空间中的既有语义结构，导致检索性能随时间下降；若通过全量重训练进行适配，又会带来较高的历史样本重编码与索引维护成本。因此，如何在不重编码历史样本的前提下实现稳定、高效的动态哈希检索，成为一个值得研究的问题。

围绕上述问题，本文从单模态图像检索、跨模态检索和系统应用三个层面展开研究。针对动态单模态图像检索任务，提出半监督概念保留哈希方法（Semi-supervised Concept Preserving Hashing，SCPH）。该方法通过为类别分配固定概念编码，在汉明空间中建立稳定语义锚点；同时引入 K近邻（K-Nearest Neighbor，KNN）伪标签机制，联合利用少量标注样本与未标注样本进行模型更新。通过综合优化样本相似性保持、目标编码一致性和哈希码均衡性，SCPH 在不重编码历史样本的条件下提升了模型对概念漂移的适应能力。

针对动态跨模态检索任务，进一步提出基于多哈希表结构的增量式跨模态哈希方法（Multi-table based Incremental Hashing，MIH）。该方法通过多表结构保留不同时期学习到的跨模态语义映射，并依据最新数据对各哈希表进行动态加权；在检索阶段，再结合查询自适应比特加权构造细粒度加权汉明距离。由此，MIH 能够在适应当前数据分布的同时保留历史语义信息，从而缓解概念漂移和灾难性遗忘带来的性能退化。

在方法研究基础上，本文面向电商场景设计并实现了一套基于动态哈希的智能检索系统，将 SCPH 与 MIH 集成到数据接入、特征提取、哈希编码、索引存储与检索服务的完整流程中，支持相似商品推荐、拍照购和自然语言搜索等功能。实验中，本文在 CIFAR-10、CIFAR-100、MIRFlickr 和 MSCOCO 数据集上构建新类别出现与数据分布变化两类概念漂移场景。结果表明，SCPH 与 MIH 在多时间步检索过程中均能保持较稳定的性能，并在概念漂移发生后表现出更好的适应能力；消融实验进一步验证了概念编码约束、多哈希表结构和查询自适应加权等关键模块的有效性。

综上，本文针对概念漂移条件下的动态哈希检索提出了两类有效方法，并通过系统实现验证了其工程可行性。相关研究为动态数据环境下构建稳定、高效的智能检索系统提供了参考。

\end{cabstract}
\ckeywords{概念漂移；在线哈希检索；半监督学习；跨模态检索；增量学习；多哈希表}
\begin{eabstract}
With the growth of Internet and mobile applications, multimedia data such as images and texts are continuously accumulated. Hashing-based retrieval is widely used in large-scale retrieval because it maps high-dimensional features into compact binary codes and supports efficient approximate nearest neighbor search in Hamming space. However, most existing methods are developed under static data assumptions and are not well suited to streaming environments where concept drift frequently occurs through newly emerging classes and changing data distributions. Once concept drift appears, the learned semantic structure becomes unstable and retrieval performance gradually degrades. Re-training from scratch is also expensive because it often requires re-encoding historical samples and rebuilding the index.

To address this problem, this thesis studies online hashing retrieval from three aspects: single-modal image retrieval, cross-modal retrieval, and system implementation. For dynamic single-modal image retrieval, a Semi-supervised Concept Preserving Hashing method (SCPH) is proposed. SCPH assigns a fixed concept code to each category to preserve stable semantic anchors in Hamming space, and introduces a K-nearest neighbor pseudo-labeling strategy to exploit unlabeled samples. By jointly optimizing similarity preservation, target-code consistency, and code balance, SCPH improves adaptation to concept drift without re-encoding historical samples.

For dynamic cross-modal retrieval, this thesis further proposes a Multi-table based Incremental Hashing method (MIH). MIH preserves cross-modal semantic mappings learned at different time steps with a multi-table structure, dynamically reweights hash tables according to their performance on newly arrived data, and introduces query-adaptive bit weighting during retrieval. In this way, MIH balances adaptation to current distributions with preservation of historical semantic information, thus mitigating concept drift and catastrophic forgetting.

Based on the proposed methods, an intelligent e-commerce retrieval system is designed and implemented. The system integrates SCPH and MIH into a complete pipeline including data ingestion, feature extraction, hash encoding, index storage, and retrieval services, and supports similar-item recommendation, image-based shopping, and natural-language search. Experiments on CIFAR-10, CIFAR-100, MIRFlickr, and MSCOCO under two concept drift settings, namely emerging classes and distribution changes, show that SCPH and MIH maintain more stable retrieval performance over multiple time steps and adapt better after drift occurs. Ablation studies further verify the effectiveness of concept-code constraints, the multi-table mechanism, and query-adaptive weighting.

In summary, this thesis proposes two effective methods for online hashing retrieval under concept drift and validates their engineering feasibility through system implementation. The study provides a practical reference for building stable and efficient intelligent retrieval systems in dynamic data environments.

\end{eabstract}
\ekeywords{Concept Drift; Online Hashing Retrieval; Semi-supervised Learning; Cross-modal Retrieval; Incremental Learning; Multi-hash Table}
